\section{The exact finite-orbit reduction of CountI}
\label{sec:count-i-finite-orbit-audit}

\status{KA1--KA4 reconstruct the full implication
\texttt{PropGK}$\Rightarrow$\texttt{CorGK}$\Rightarrow$\texttt{CountI}.
The arithmetic construction of \texttt{PropGK} remains an input.
In particular, the assertions below do not supply \texttt{OneChar},
the rank-three congruence-number theorem, or the automorphy and
rigidity hypotheses needed to construct them.  This isolates the
deep producer from the finite aggregation at the end of the CountI
branch.  Every coefficient ring and finite set of primes in the
original predicate remains arbitrary.  KA5--KA6 prove that Case III
is empty for targets of power-of-two order, including the principal
target, and carry that specialization through the actual height and
cohomological-window consumers.}

\subsection{Interfaces and quantifier order}

Use the original host notation of HV1: a fixed primitive character
$\chi$, conductor $q$, order $o$, exponent $s\ge2$, and auxiliary
primes $q_1,d$.  Write
\[
 N=\operatorname{lcm}(2^{s+2},q)q_1d,
 \qquad G_N=(\mathbb Z/N\mathbb Z)^\times.
 \tag{KA.1}
\]
The host subgroup $H_h\le G_N$ is cut out by the sign condition at
$2^{s+2}$, the kernel of $\chi$, and the square conditions at $q_1,d$.
For any source coefficient DVR $O$ at an odd prime $\ell$, let
$X_h(O,e)$ denote the characters trivial on $H_h$ and satisfying
$\psi(-1)=(-1)^e$.  Their original extended Selmer length and
normalized length are denoted by $L_{\ell,O}(\psi,e)$ and
$v_{\ell,O}(\psi,e)$, exactly as in (HV.3)--(HV.5).  Before
finiteness is proved, only the extended length is used.  The
normalized quantity in the source is always nonnegative.

For $J\le G_N$, put
\[
 \mathcal O_O(J)=\{\psi:G_N\longrightarrow O^\times:
                         \ker\psi=J\}.
 \tag{KA.2}
\]
This is the literal source \texttt{orb}.  It can be empty; no common
coefficient field containing all roots of unity is assumed.

For a number field $K$, write $D_K=|\operatorname{disc}K|$.
The predicate \texttt{GD} consists of the following conditions:
$\ell$ is odd; $e\ge1$; $N$ is even; $\ell\nmid\operatorname{ord}\psi$;
the $\ell$-primary CRT component of $\psi$ has square one;
$\psi(-1)=(-1)^e$; $K$ is imaginary quadratic;
$\ell\nmid D_Kh_K$; $(D_K,N)=1$;
\[
 \ker\overline\chi_{\mathrm{cyc},\ell}
   \cap\ker\operatorname{dir}(\chi_P(\ell,\psi))
       \not\subseteq G_K;
 \qquad \ell-1\nmid2e.
 \tag{KA.3}
\]
Here $\chi_P$ is the source character obtained by deleting the
$\ell$-power-order part and the local $\ell$-component.  In
particular its conductor divides $N$.

The exact useful form of \texttt{CorGK} is this.  For every fixed
number field $K$ there are integers $c_5,A\ge0$, chosen before
$\Omega,e,N,J,p,p'$, the coefficient rings, and the finite prime
set $\mathcal L$, such that if $p\ne p'$ are primes not dividing
$30ND_K$ and for each $\ell\in\mathcal L$ at least one character in
$\mathcal O_{O_\ell}(J)$ satisfies \texttt{GD}, then all lengths
in these orbits are finite and
\[
 \sum_{\ell\in\mathcal L}\log\ell
       \sum_{\psi\in\mathcal O_{O_\ell}(J)}v_{\ell,O_\ell}(\psi,e)
 \le 2c_5\varphi([G_N:J])\,[N(e+1)]^A
                      (\log p+\log p').
 \tag{KA.4}
\]
Only the existence of one admissible orbit member at each prime is
required to apply this exact input; its conclusion covers the whole
orbit in that coefficient ring.

The stronger input \texttt{PropGK} supplies, for every such
$K,e,N,J$ and every prime $p\nmid30ND_K$, a nonzero integer $b_p$,
chosen before every later $\ell\ne p$ and coefficient ring, with
\[
 \log|b_p|\le c_5\varphi([G_N:J])\,[N(e+1)]^A\log p,
 \qquad
 \sum_{\psi\in\mathcal O_{O_\ell}(J)}v_{\ell,O_\ell}(\psi,e)
       \le 2\operatorname{ord}_{\ell}(b_p),
 \tag{KA.5}
\]
and finiteness of those lengths whenever the same \texttt{GD}
antecedent holds.  The second assertion is conditional on
$\ell\ne p$; this exception is why two auxiliary primes occur in
(KA.4).  The source \texttt{CorG}, with fixed exponents and a
different constant, is not substituted for \texttt{CorGK}.

\paragraph{KA1: total mass of character-kernel orbits.}
\label{claim:KA1}
Let $G$ be a finite abelian group, and let $\mathcal T$ be any
finite set of subgroups $J\le G$ such that $G/J$ is cyclic.  Then
\[
       \sum_{J\in\mathcal T}\varphi([G:J])\le |G|.
 \tag{KA.6}
\]
If $J$ is the kernel of a character with values in the units of a
characteristic-zero domain, then $G/J$ is cyclic.

\begin{proof}
A finite subgroup of the multiplicative group of a field is cyclic.
One proof is to write its exponent as $m$.  If the group had more
than $m$ elements, the polynomial $X^m-1$ would have too many roots;
the structure theorem for finite abelian groups then shows that a
group with cardinality equal to its exponent is cyclic.  Apply this
to the image of the character, using the fraction field of the
coefficient domain.

For a cyclic quotient $G/J$ of order $m=[G:J]$, exactly
$\varphi(m)$ complex characters are faithful: a generator may be
sent to any primitive $m$th root of unity.  Composition with the
quotient map identifies these with the complex characters of $G$
whose kernel is exactly $J$.  For different $J$ these sets are
disjoint.  A finite abelian group has exactly $|G|$ complex
characters, as follows by decomposing $G$ into cyclic factors and
choosing the image of one generator in each factor.  Counting the
disjoint subsets gives (KA.6).  This argument uses complex
characters only as a counting device; it does not enlarge or
identify the original coefficient DVRs.
\end{proof}

\paragraph{KA2: the exact Case-I host supplies GD.}
\label{claim:KA2}
Fix $\chi$ and let $p_K=p_K(\chi)$, the source prime satisfying
$p_K\equiv3\pmod4$, $p_K\ge7$, and $p_K\nmid q$.
Let $K_\chi=K_*(p_K)$ be the source splitting-field model of
$X^2+p_K$, and let $h_*=h(K_\chi)$.
For every original host $h$ with character $\chi$ and
$\texttt{Big}(t)$, every odd prime $\ell$ in Case I, every
$1\le e<t$, every original coefficient ring $O$, and every
$\psi\in X_h(O,e)$, the datum
$(\Omega,K_\chi,\ell,e,\psi)$ satisfies the literal \texttt{GD}.
The statement is valid for every embedding of $K_\chi$ into the
given algebraic closure $\Omega$.

\begin{proof}
There are infinitely many primes congruent to $3$ modulo $4$:
if a finite list included all of them, a prime factor congruent to
$3$ modulo $4$ of four times their product minus one would be a
new one.  Including $3$ and the prime divisors of $q$ among the
excluded finite list supplies an eligible $p_K$.  Thus the
source minimum is well defined.  Since $-p_K$ is squarefree and
congruent to $1$ modulo $4$, the ring of integers of
$\mathbb Q(\sqrt{-p_K})$ has basis
$1,(1+\sqrt{-p_K})/2$ and discriminant $-p_K$.
The two roots $\pm\sqrt{-p_K}$ show that its splitting field is
the literal $K_*(p_K)$ up to a rational-field
isomorphism.  Hence $K_\chi$ is imaginary quadratic and
$D_{K_\chi}=p_K$.

The source exceptional set is
\[
 \mathcal B_\chi=\{3\}\cup
    \{\text{odd prime divisors of }oqp_Kh_*p_1\}.
 \tag{KA.7}
\]
Case I means $\ell\notin\mathcal B_\chi$ and
$\ell-1\nmid2e$.  Therefore
$\ell\nmid oqp_Kh_*$ and $\ell\ne p_K$.
The assumptions $\texttt{Big}(t)$ give $q_1,d>p_K$;
using $p_K\nmid q$ and (KA.1), we obtain $p_K\nmid N$.
These facts prove the two discriminant/class-number conditions
in \texttt{GD}.

The host quotient $G_N/H_h$ injects into the product of the
quotient of $(\mathbb Z/2^{s+2})^\times$ by $\langle-1\rangle$,
the image of $\chi$, and the two order-two square quotients.
Their orders are $2^s,o,2,2$, respectively.  Thus every
character trivial on $H_h$ has order dividing $2^{s+2}o$;
the odd prime $\ell\nmid o$ cannot divide its order.
Since $\ell\nmid q$, its CRT factor in the host quotient is
trivial unless $\ell=q_1$ or $\ell=d$.  In either exceptional
possibility the quotient is a square quotient, of exponent two.
Thus the local $\ell$-component of $\psi$ has square one.
Also $16\mid N$, and membership in $X_h(O,e)$ supplies the
required parity.  The oddness of $\ell$, positivity of $e$,
and nondivisibility $\ell-1\nmid2e$ are already hypotheses.

It remains to prove the noncontainment in (KA.3), which cannot be
replaced by the assertion that the two characters are nontrivial.
At a place above $p_K$, the quadratic extension $K_\chi/\mathbb Q$
is ramified.  For example, $X^2+p_K$ is Eisenstein over
$\mathbb Q_{p_K}$ and its root has valuation $1/2$; the extension
is totally ramified.  Consequently the absolute inertia group
at this place acts nontrivially on $K_\chi$.
The mod-$\ell$ cyclotomic character is unramified there because
$p_K\ne\ell$.  The Dirichlet character $\chi_P(\ell,\psi)$ is
also unramified there because its conductor divides $N$ and
$p_K\nmid N$.  This inertia group lies in both kernels in
(KA.3), while one of its elements is outside $G_{K_\chi}$.
This proves the noncontainment.  All assertions concern the
given embedded fields and their decomposition/inertia groups;
conjugating the chosen embedding preserves them.
\end{proof}

\paragraph{KA3: one integer at each of two primes gives CorGK.}
\label{claim:KA3}
The literal \texttt{PropGK} implies the literal \texttt{CorGK},
with the same constants $c_5,A$.

\begin{proof}
If $\mathcal L$ is empty, both the finiteness assertion and the
sum bound are immediate; the right side of (KA.4) is nonnegative.
Otherwise the assumed \texttt{GD} witness for one prime implies
$e\ge1$, so \texttt{PropGK} applies.  Use it at $p$ and $p'$,
obtaining $b_p,b_{p'}\ne0$ before inspecting any $\ell$ or $O_\ell$.
For each $\ell\in\mathcal L$, at least one of $\ell\ne p$ and
$\ell\ne p'$ holds.  Applying the corresponding integer bound
proves finiteness and, by nonnegativity of both valuations,
\[
  \sum_{\psi\in\mathcal O_{O_\ell}(J)}v_{\ell,O_\ell}(\psi,e)
    \le2\operatorname{ord}_\ell(b_p)
                      +2\operatorname{ord}_\ell(b_{p'}).
 \tag{KA.8}
\]
For every nonzero integer $b$ and finite set of distinct primes,
\[
 \prod_{\ell\in\mathcal L}\ell^{\operatorname{ord}_\ell(b)}
       \mid |b|,
 \qquad
 \sum_{\ell\in\mathcal L}\operatorname{ord}_\ell(b)\log\ell
       \le\log|b|.
 \tag{KA.9}
\]
The divisibility follows from unique factorization and the
logarithmic inequality follows by taking logarithms of positive
integers.  Multiplying (KA.8) by $\log\ell\ge0$, summing, and
applying (KA.9) to both integers gives
$2\log|b_p|+2\log|b_{p'}|$.  The first inequalities of (KA.5)
are exactly (KA.4).  This also shows why the same two integers
must work for all later coefficient rings and primes: integers
chosen after $\ell$ could not be combined using (KA.9).
\end{proof}

\paragraph{KA4: complete original-host aggregation to CountI.}
\label{claim:KA4}
The literal \texttt{CorGK} implies the literal \texttt{CountI}.
Hence \texttt{PropGK} implies \texttt{CountI} as well.

\begin{proof}
Fix the primitive target $\chi$, and apply \texttt{CorGK} to
the fixed field $K_\chi=K_*(p_K(\chi))$.  Obtain $c_5,A\in\mathbb N$.
Choose distinct primes
\[
       r_1>30qp_K,\qquad r_2>r_1.
 \tag{KA.10}
\]
Euclid's theorem supplies them.  These choices depend only on
$\chi$.  Put $B=2c_5(\log r_1+\log r_2)\ge0$ and choose
$k\in\mathbb N$ so that $B\le16^k$.  Finally put
\[
                    C=k+1+A.
 \tag{KA.11}
\]
This has made every required existential choice before the host,
the cutoff $t$, the twist $e$, the algebraic closure, the family of
coefficient rings, and the finite set of primes.

Now take any original host of character $\chi$ with
$q_1,d>r_1,r_2$, any $\texttt{Big}(t)$, any integer
$1\le e<t$, any algebraic closure $\Omega$, any coefficient
family $(O_\ell)_\ell$, and any finite set $\mathcal L$ of odd
Case-I primes.  Embed $K_\chi$ into $\Omega$; the following
estimates are independent of that choice by KA2.
Every prime dividing $N$ is either $2$, a divisor of $q$,
$q_1$, or $d$.  Since each $r_i$ exceeds $30qp_K$ but is smaller
than $q_1,d$, neither $r_i$ divides $30Np_K$.  Thus the same two
primes are admissible in (KA.4) for every such host.

Let $\mathcal T$ be the union, over $\ell\in\mathcal L$, of
the kernels of all characters in $X_h(O_\ell,e)$.  For
$J\in\mathcal T$, let
\[
 \mathcal L_J=\{\ell\in\mathcal L:
     \text{some }\psi\in X_h(O_\ell,e)\text{ has kernel }J\}.
 \tag{KA.12}
\]
For each $\ell\in\mathcal L_J$ choose such a character only to
verify the antecedent of \texttt{CorGK}; KA2 supplies its
\texttt{GD}.  Applying (KA.4) to $J$ and $\mathcal L_J$ gives
finiteness for every character in its full coefficient-ring orbit
and the estimate
\[
 \sum_{\ell\in\mathcal L_J}\log\ell
       \sum_{\psi\in\mathcal O_{O_\ell}(J)}v_{\ell,O_\ell}(\psi,e)
       \le B\,\varphi([G_N:J])\,[N(e+1)]^A.
 \tag{KA.13}
\]
Every host character lies in precisely one kernel fiber, which is
a subset of the corresponding full orbit.  The quantities being
summed are nonnegative; therefore enlarging each fiber to that
orbit and interchanging two finite sums is valid.  It gives
\begin{align*}
 &\sum_{\ell\in\mathcal L}\log\ell
                   \sum_{\psi\in X_h(O_\ell,e)}v_{\ell,O_\ell}(\psi,e)\\
 &\quad\le\sum_{J\in\mathcal T}
       \sum_{\ell\in\mathcal L_J}\log\ell
                   \sum_{\psi\in\mathcal O_{O_\ell}(J)}v_{\ell,O_\ell}(\psi,e)\\
 &\quad\le B[N(e+1)]^A
                    \sum_{J\in\mathcal T}\varphi([G_N:J])
       \le BN[N(e+1)]^A.
 \tag{KA.14}
\end{align*}
The last step is KA1 and $|G_N|=\varphi(N)\le N$.
Notice that $\mathcal T$ is a set of subgroups of the one fixed
group $G_N$; no identification between the coefficient rings is
used.  Finiteness for all original host characters was already
obtained from (KA.13).

Finally $16\mid N$ gives $N\ge16$, and $e<t$ with integer
$e,t$ gives $e+1\le t$.  Hence
\[
 BN[N(e+1)]^A
 \le N^kN(Nt)^A
 \le (Nt)^{k+1+A}=(Nt)^C.
 \tag{KA.15}
\]
This is exactly the numerical conclusion of \texttt{CountI},
together with its finiteness conclusion and the required order
of all quantifiers.  KA3 gives the stated further implication
from \texttt{PropGK}.
\end{proof}

\paragraph{KA5: Case III vanishes for every target of power-of-two order.}
\label{claim:KA5}
Suppose $\operatorname{ord}(\chi)$ is a power of two, including
$\operatorname{ord}(\chi)=1$.  Then, on every original host of
character $\chi$, there is no $\psi\in X_h(O,e)$ satisfying
\texttt{CaseIII}, for any coefficient ring, prime or twist.
Consequently the fixed-$\chi$ specialization of \texttt{CountIII}
holds with exponent zero.

\begin{proof}
The quotient calculation in KA2 shows
\[
  \operatorname{ord}(\psi)\mid2^{s+2}\operatorname{ord}(\chi).
 \tag{KA.16}
\]
(The source \texttt{hostComp} proves the stronger exponent
$2^{s+1}\operatorname{ord}(\chi)$, which is not needed.)
If the right side is a power of two, so is the positive integer
$\operatorname{ord}(\psi)$.  The literal \texttt{CaseIII}
requires this order to be odd and requires $\psi\ne1$.
An odd divisor of a power of two is one, whereas a character of
order one is the trivial character.  This is a contradiction.
The numerical and finiteness assertions in the fixed-target
\texttt{CountShape} are therefore implications with false
antecedent.  Choosing $C=0$ satisfies all remaining quantifiers.
No assertion about \texttt{CountIII} for targets with nontrivial
odd-order part has been made.
\end{proof}

\paragraph{KA6: the specialized height/window cut.}
\label{claim:KA6}
For any fixed target $\chi$ of power-of-two order, the following
are sufficient for the original \texttt{PointHV} assertion at
that target:
\[
 \operatorname{CorGK}_{K_\chi},\qquad
 \texttt{PMExAlt},\qquad \operatorname{CountV}_{\chi}.
 \tag{KA.17}
\]
Here the subscript means the exact original predicate with only
its outer field or target variable fixed, leaving every later
quantifier unchanged.  One can replace the first premise by
$\operatorname{PropGK}_{K_\chi}$.  With the supplied TW/OB/VW
inputs and their unchanged prime-pair and window hypotheses, the
same premises supply $\texttt{VWinAt}(\chi,401)$.  Thus no
Case-III seed or promodularity theorem is needed on this selected
cohomological route to the principal target.

\begin{proof}
KA4 uses only \texttt{CorGK} at the single field $K_\chi$, and
therefore supplies $\operatorname{CountI}_\chi$ from that
specialization.  YX2 supplies the literal \texttt{CountII} from
\texttt{PMExAlt}, so in particular its $\chi$ specialization.
KA5 supplies $\operatorname{CountIII}_\chi$ with exponent zero.
The previously reconstructed \texttt{CountZ}, \texttt{CountVs}
and \texttt{CountIV} remain available in their original scopes.

To check that fixed-target counts really suffice, inspect the
quantifiers of the actual consumers, not merely their universal
wrapper.  The proof of \texttt{sumOut} fixes $\chi$ and calls
\texttt{CountI}, \texttt{CountZ}, \texttt{CountII} and
\texttt{CountIII} only at this same $\chi$; it then performs
the finite prime, twist and character sums of HV9.  The proof of
\texttt{sumIn} likewise uses \texttt{CountII}, \texttt{CountV},
\texttt{CountVs} and the layer count only at this $\chi$.
In \texttt{point}, after fixing $(\chi,j)$, all exceptional
constants are chosen from \texttt{sumIn} at
$(\chi,\ell,j(\ell))$ and the outside constant is chosen from
\texttt{sumOut} at $\chi$.  Thus the proofs of HV8--HV10 run
unchanged with these specialized predicates.  They give precisely
the same field-model, coefficient, layer-function and uniform
host-threshold quantifiers of $\operatorname{PointHV}_\chi$.

Finally \texttt{thmHV} uses \texttt{PointHV} only at the supplied
target $\chi$, with
$j(\ell)=\lfloor\log_\ell R_\ell(L_\chi)\rfloor+3$ in the
original integer-logarithm convention.  The source \texttt{jointWin\_of}
calls this height estimate at the same target, and
\texttt{vwinAt\_of\_joint} uses the joint-window hypothesis only
at that same target.  The already supplied two-primary and
odd-minus inputs are global and may of course be specialized.
Consequently the full field/window chain remains valid at
$\chi$, including the original real continuum of scales and
the cutoff exponent $400$ when $C_t=401$.  None of these
arguments provides the still-separate common arithmetic family,
gap, theta, denominator or metric/spectral branches required
by the final theorem.
\end{proof}

\paragraph{Why this does not remove cases I, II or V for RH.}
The principal target has $p_K=7$, $p_1=3$, and
$h(\mathbb Q(\sqrt{-7}))=1$, hence $\mathcal B_1=\{3,7\}$.
The class-number identity has the elementary reduced-form check:
if $[a,b,c]$ is a reduced primitive positive binary quadratic form
of discriminant $-7$, then $|b|\le a\le c$ implies
$3a^2\le7$.  Thus $a=1$, $b=\pm1$, and $c=2$; the usual reduced
boundary convention selects $[1,1,2]$.  The standard
ideal-class/binary-form correspondence therefore gives one class.
Thus Case V specializes to $\ell\in\{3,7\}$ and $e=1$;
its dyadic conductor is still unbounded.  For a concrete
nonempty instance take $q_1=11$, $d=23$, $t=2$ and host exponent
$s=12$.  A primitive even character modulo $2^{14}$, multiplied
by the Legendre character modulo $11$ and the trivial character
modulo $23$, is an allowed odd host character.  Choose an
original $3$-adic coefficient ring containing its roots of unity.
For $\ell=3$ its source invariants satisfy
\[
  f_0=2,\qquad n_0=3,\qquad c_2=14,\qquad
  R_1=16\cdot3^2\cdot4^3=9216<2^{14}.
 \tag{KA.18}
\]
It is therefore in the literal \texttt{CountV} case.  The host
is \texttt{Big}(2), since $5,7,3<11,23$ and $q=o=h_*=1$.
A primitive even character of conductor $2^{14}$ exists by
$(\mathbb Z/2^{14})^\times\cong C_2\times C_{2^{12}}$:
send $-1$ to $1$ and $5$ to a primitive $2^{12}$th root.

Case I is nonempty already with $e=1$, $\ell=5$ and the odd
Legendre character modulo $11$, since $5\notin\{3,7\}$ and
$4\nmid2$.  Case II is nonempty with $e=2$, $t=3$,
$\ell=3$, and the product of the Legendre characters modulo
$11$ and $23$; its order is two and its parity is even.
These examples only prevent an erroneous empty-case elimination;
they assert no cohomological lower bound.  In particular the
principal route still encounters characteristic three, so a
support theorem restricted to $\ell>3$ cannot by itself supply
its whole modular-support obligation.

\subsection{What this closes and what it leaves}

For a general target, the remaining arithmetic input on the
Case-I branch may now be named precisely:
\texttt{PropGK} for the original coefficient and host interfaces.
An even narrower sufficient input for a fixed target is its
specialization to the single quadratic field $K_*(p_K(\chi))$;
the argument above never varies $K$ after choosing $\chi$.
For a proof of every Dirichlet target, the field specialization
must be supplied for every $\chi$.

The source constructs \texttt{PropGK} from \texttt{OneChar} and
\texttt{CongruenceNumberRankThree}.  That producer controls one
nonzero integer simultaneously for all later coefficient primes
and contains the rank-two/rank-three automorphic comparison.
It is not a consequence of the finite orbit count proved here.
The present reduction also does not identify rank-four RX/CX
statements with rank-three Kato or host-frame statements.
For ordinary RH and, more generally, power-of-two-order targets,
KA5--KA6 remove the separate Case-III producer from the selected
cohomological route.  They do not remove it from unrestricted
Dirichlet GRH, and they do not assert that every other arithmetic
branch is complete.

\paragraph{Exact source record.}
All source references below are at commit
\texttt{7cbf7a0c32fca1badf05931879a37c3f8a2b03b4} of
\texttt{CocoSgt/riemann-hypothesis-lean}.  KA1 follows
\texttt{HeightAndVanishing/OrbSum.lean:orbSum} and
\texttt{OrbFacts.lean:orbFacts}.  KA2 reconstructs
\texttt{GdOfHost.lean:gdOfHost}, with the actual host definitions
in \texttt{DirichletPrimeToLPartMulLPowerPart.lean},
\texttt{BchiFacts.lean:bchiFacts},
\texttt{HostIndex.lean:hostIndex}, and
\texttt{SplittingFieldIsoQuadratic.lean:kstar}.
KA3 reconstructs \texttt{UnitaryThreeGroups/CorGSumLog.lean:corGK}.
KA4 reconstructs \texttt{UnitaryThreeGroups/CountI.lean:countI}.
The exact predicates are \texttt{UnitaryThreeGroups/KatoUse.lean}
and \texttt{HeightAndVanishing/Counts.lean}.
The surviving producer is
\texttt{UnitaryThreeGroups/OrbEquiv.lean:propGK}, with the
actual \texttt{OneChar} supplier under
\texttt{ModularityArgument/KatoQuadraticCharacters/OneChar.lean}.
The accompanying provenance file records exact line ranges,
file hashes, Git blob identifiers and links for each claim.
KA5 additionally uses \texttt{HostComp.lean:hostComp} and the
literal \texttt{Host.CaseIII} definition.  KA6 is checked against
\texttt{SumOut.lean:sumOut}, \texttt{FinalLe.lean:sumIn},
\texttt{PointArith.lean:point}, \texttt{ThmHV.lean:thmHV},
\texttt{PrimeSetBoundDataOfPrimitiveChar.lean:jointWin\_of}, and
\texttt{CMFieldExponentialBounds/AdmissibleOfPageN.lean:vwinAt\_of\_joint}.
